\documentclass[11pt,letterpaper]{article}
\usepackage[T1]{fontenc}
\usepackage{lmodern,microtype}
\usepackage[margin=1in]{geometry}
\usepackage{amsmath,amssymb,amsthm,mathtools}
\usepackage{booktabs,tabularx,longtable,array}
\usepackage{enumitem,xurl,hyperref,fancyhdr}
\hypersetup{hidelinks,pdftitle={Singular-Safe Port Compression and Finite-State Multiplicity Registers},pdfauthor={Research draft prepared with OpenAI ChatGPT for ProveIt}}
\pagestyle{fancy}\fancyhf{}
\fancyhead[L]{\small ProveIt / Unknot recognition}
\fancyhead[R]{\small Port-register compression}
\fancyfoot[C]{\thepage}
\setlength{\headheight}{14pt}
\setlength{\emergencystretch}{2em}
\setcounter{tocdepth}{1}
\allowdisplaybreaks[2]
\newcommand{\F}{{\mathbb F_2}}
\newcommand{\N}{\mathbb N}
\newcommand{\rk}{\operatorname{rank}}
\newcommand{\im}{\operatorname{im}}
\newcommand{\Span}{\operatorname{span}}
\newcommand{\one}{\mathbf 1}
\newcommand{\id}{\operatorname{id}}
\newcommand{\Kh}{\operatorname{Kh}}
\newcommand{\code}[1]{\texttt{\detokenize{#1}}}
\newtheorem{theorem}{Theorem}[section]
\newtheorem{lemma}[theorem]{Lemma}
\newtheorem{proposition}[theorem]{Proposition}
\newtheorem{corollary}[theorem]{Corollary}
\theoremstyle{definition}
\newtheorem{definition}[theorem]{Definition}
\newtheorem{example}[theorem]{Example}
\theoremstyle{remark}
\newtheorem{remark}[theorem]{Remark}
\numberwithin{equation}{section}
\title{\LARGE Singular-Safe Port Compression\\ and Finite-State Multiplicity Registers\\[0.5em]
\large Exact homology beyond direct-sum sharing, with a conditional\\ route to quasi-polynomial unknot recognition}
\author{Research continuation for the ProveIt project}
\date{8 October 2026\quad /\quad Reproducible research draft}
\begin{document}
\maketitle
\begin{abstract}
Explicit chain-object multiplicity remains an asymptotic obstruction in the
ProveIt Khovanov scanner. Literal direct-sum sharing does not describe coupling
between copies. We develop an exact scalar-complex backend for a different
representation: small template differentials repeated over exponentially many
Boolean words, connected by a bounded number of rank-one channels, called
ports. The port coordinates are given by layered weighted automata over
$\F$, optionally restricted to an explicitly supplied deterministic language.

A singular-safe bordered-rank formula reduces the calculation to a matrix
with at most $2r$ rows and columns, where $r$ is the port count. Its use does
not require an inverse of the port interaction matrix. Finite-state row
spaces and parity pairings construct this matrix without enumerating the
multiplicity domain. The same primitives check the full square-zero identity
and minimize the port count exactly for the chosen base differential,
without increasing register width. We prove a sharp rank-budget threshold
for decision saturation and a tensor-amplification obstruction to automatic
closure of the representation.

The resulting algorithm is polynomial in the explicit register presentation,
not in the virtual chain dimension. A certified quasi-polynomial-time
presentation producer would therefore yield a quasi-polynomial recognizer;
no such producer for arbitrary diagrams is established here. The package
contains proofs, standard-library Python, 48 unit tests, 4,352 exhaustive
rank-update comparisons, 404 braid-closure cube audits, and both favorable
and unfavorable paired measurements. Exponentially large abstract examples
are not presented as knot benchmarks or as evidence of an unrestricted
recognition bound.
\end{abstract}
\clearpage
\tableofcontents
\bigskip
\noindent\textbf{Reading guide.}
Sections~\ref{sec:border}--\ref{sec:complexity} give the exact algebra,
finite-state construction, validity test, and presentation-size bound.
Sections~\ref{sec:budget} and~\ref{sec:amplification} establish the sharp
cancellation budget and a limitation on propagating scalar ports.
For implementation and evidence, read Sections~\ref{sec:integration}
and~\ref{sec:experiments}, then Appendix~\ref{app:files}.
The constructive diagram-to-presentation question is isolated in
Theorem~\ref{thm:conditional} and the research agenda.
\newpage

\section{The bottleneck and the precise advance}
\label{sec:advance}

\subsection{What was inspected}
The inspected ProveIt implementation stores the objects of each scanning
stage in explicit lists, interns matching types, and uses sparse, prioritized
unit cancellation. The complete default method still has an exponential
general guarantee. Our source audit covers the project README, the beginning
of \code{scan_fast.py}, and the maintained structural and radical-transfer
sections; it is not a full checkout audit \cite{repo-readme,repo-scan,
repo-structural,repo-radical}. The observed branch commit was
\code{81a387aada56bb180649a98c21c078cb89efcfde}. The scanner blob was checked at
that commit. Other inspected blobs and the exact scope of their retrieval
are recorded separately in \code{provenance.json}.

Two existing ideas are particularly relevant. Component sharing compresses
\emph{literal direct sums}: equal differential components may be represented
once with a multiplicity. Residue prediction determines how many matching
copies survive scalar-unit cancellation, but not all surviving maps. The
present work changes the multiplicity representation rather than merely
accelerating another coefficient operation. It permits nonzero maps between
copies, including examples whose entire differential support graph is
connected.

This is a different contribution from the earlier twist, rank-two braid,
extremal-window, graded-occupancy, and common-disk radical-transfer routes.
It also has a different limitation: our verified algorithm operates on
closed, scalar chain complexes. A closed homology rank is not a replacement
for the relative chain-homotopy data of an open tangle.

\subsection{Relation to the literature}
Bar--Natan's local simplification method provides the computational setting
\cite{barnatan}. Kelom\"aki and Sch\"utz show that the usual scanning and
explicit-basis divide-and-conquer algorithms can require exponential time
even on certain positive or alternating three-braids, while a different
algorithm computes three-braid homology in polynomial time \cite{kelomaki}.
The distinction between an explicit basis and a succinct description is
therefore substantive, not merely a programming preference.

The linear algebra below uses classical rank-factorization and bordered
matrix operations. The register representation is a layered version of a
weighted automaton; forward-space methods have established polynomial
algorithms \cite{kiefer}, and matrix-product representations also fit the
tensor-train viewpoint \cite{oseledets}. We do not claim these ingredients,
Lucas parity, or the elementary rank inequality as new discoveries. The
contribution is their precise synthesis into a singular-safe, exact,
proof-audited interface for coupled multiplicities, including optimal
fixed-base port reduction, domain restriction, and explicit failure modes.
Literature priority for this combined formulation has not been established.

Lackenby's July 2026 hierarchy paper is a separate route. Section 9 explicitly
leaves several operation costs and controlling parameters outside a complete
running-time estimate \cite{lackenby}. The scalar-complex results here neither
supply those missing geometric operations nor contradict the announced
quasi-polynomial program motivating the repository.

\subsection{Status of the claims}
\begin{center}
\begin{tabularx}{\linewidth}{@{}p{0.25\linewidth}Xp{0.19\linewidth}@{}}
\toprule
Claim & Scope & Status\\
\midrule
Small bordered core & Arbitrary binary low-rank matrix updates, including singular interaction & Proved and tested\\
Succinct homology & Valid port-register complexes with explicit finite-state descriptions & Proved and implemented\\
Minimum port count & Fixed base differential and homogeneous scalar ports & Proved and implemented\\
Knot interpretation & A separate certified association with closed Khovanov homology & Conditional on provenance\\
General quasi-polynomial recognition & Arbitrary ordinary knot diagrams & Not established by this work\\
\bottomrule
\end{tabularx}
\end{center}

The measurements support a potential kernel improvement, not a default-backend
change. A producer that first expands a small cube and then applies our
interface is much slower than direct rank computation on that cube. The
useful regime is one in which the producer avoids expansion in the first
place.

\section{Scalar complexes, templates, and ports}
\label{sec:model}

All coefficient arithmetic is over $\F$. We use cochain conventions:
a differential raises homological degree by one. Quantum grading is not
needed for the total-rank decision studied here and is not computed by the
prototype.

\begin{definition}[Port-register presentation]
For each block $\nu$, choose a finite graded space $T_\nu$ of dimension
$d_\nu$, a degree-one matrix $A_\nu$ satisfying $A_\nu^2=0$, and a domain
$\Omega_\nu\subseteq\{0,1\}^{m_\nu}$. In the basic model the domain is all
words. The language-restricted model supplies an explicit layered
deterministic automaton accepting exactly $\Omega_\nu$.

Set
\begin{equation}
 C=\bigoplus_\nu T_\nu\otimes\F^{\Omega_\nu},\qquad
 A=\bigoplus_\nu A_\nu\otimes I_{|\Omega_\nu|},\qquad D=A+UV,
 \label{eq:model}
\end{equation}
where $U$ has $r$ columns and $V$ has $r$ rows. Port $j$ has a prescribed
integer $\ell_j$: its row in $V$ is supported in degree $\ell_j$, and its
column in $U$ is supported in degree $\ell_j+1$. At each block and word,
a weighted register simultaneously specifies the $d_\nu r$ coordinates
of $U$ and the $d_\nu r$ coordinates of $V$. Validity requires $D^2=0$.
\end{definition}

Thus a port is a homogeneous rank-one channel, not a topological boundary
circle or a cobordism generator. Several ports can be linearly dependent,
and cancellation can make $UV=0$ even when neither factor is zero.
No positivity is assumed for differential entries. Multiplicity counts,
by contrast, are ordinary nonnegative integers and are never reduced
modulo two.

Write
\begin{equation}
 N=\sum_\nu d_\nu|\Omega_\nu|,\quad
 q_\nu=\rk A_\nu,\quad
 q=\sum_\nu q_\nu|\Omega_\nu|,\quad
 \beta_0=\sum_\nu(d_\nu-2q_\nu)|\Omega_\nu|.
 \label{eq:basecounts}
\end{equation}
The potentially enormous integers $N,q,\beta_0$ have small binary
representations relative to the explicitly listed register lengths.
The main target is
\begin{equation}
 \beta(D)=\dim_\F H(C,D)=N-2\rk D.
 \label{eq:beta}
\end{equation}
The last equality uses $D^2=0$, hence $\im D\subseteq\ker D$.
Applying it to an arbitrary matrix without the square-zero hypothesis
would not compute homology.

\subsection{The knot criterion and its provenance requirement}
For a validated one-component classical knot diagram, total unreduced
Khovanov rank two over $\F$ detects the unknot. Here is the precise chain
of implications used by the project. Shumakovitch's splitting over
characteristic two gives twice the reduced total rank
\cite[Corollary 3.2.C]{shumakovitch}. If the reduced rank over $\F$ is one,
the universal coefficient theorem bounds its rational rank by one.
That rational rank is nonzero: the reduced Jones Euler characteristic
has value one at its standard identity specialization for a knot.
Kronheimer--Mrowka's detection theorem then implies that the knot is
trivial \cite{kronheimer}. The unknot itself has unreduced rank two.

For completeness, the characteristic-two splitting comes from the marked
point map $X$ and the dot-derivative map $\nu$, which commute with the chain
differential and satisfy
\[
 X^2=\nu^2=0,\qquad X\nu+\nu X=I.
\]
These identities split the complex into two isomorphic copies of the
reduced complex, with quantum shifts ignored. They explain why the factor
two is a mathematical feature, not a numerical convention.

\begin{remark}[An abstract rank-two complex is not a knot certificate]
The presentation \eqref{eq:model} alone says nothing about which diagram,
if any, it describes. The implementation deliberately emits
\code{topological_verdict: null}. A trusted diagram-to-complex constructor,
or a separately verified equivalence certificate, is needed before the
rank criterion can be used as a topological verdict. Link inputs also
require their own interpretation; the one-component promise is not dropped.
\end{remark}

\section{A bordered-rank formula without an inverse assumption}
\label{sec:border}

\begin{lemma}[Reflexive inverse from a rank minor]
\label{lem:inverse}
Let $A\in\F^{a\times b}$ have rank $q$. There is a matrix
$h\in\F^{b\times a}$ satisfying
\begin{equation}
 AhA=A,\qquad hAh=h.
 \label{eq:reflexive}
\end{equation}
It is constructible by elementary elimination. If $I,J$ select an
invertible $q\times q$ rank minor $M=A[I,J]$, then
$h=E_JM^{-1}E_I^T$ is such a matrix.
\end{lemma}
\begin{proof}
The selected columns span the column space of $A$, and the selected rows
of those columns are independent. Consequently
$A=A[:,J]M^{-1}A[I,:]$. Substitution gives $AhA=A$. The second identity
follows by replacing the middle $E_I^TAE_J$ with $M$. The empty rank-zero
minor gives $h=0$, including zero-dimensional cases. This is a reflexive
generalized inverse, not a positive-definite or Moore--Penrose construction.
\end{proof}

\begin{theorem}[Singular-safe port core]
\label{thm:border}
Let $A\in\F^{a\times b}$, $U\in\F^{a\times r}$, and
$V\in\F^{r\times b}$. Choose $h$ as in Lemma~\ref{lem:inverse}, and define
\begin{equation}
 U_0=(I-Ah)U,\qquad V_0=V(I-hA),\qquad E=I_r+VhU.
 \label{eq:residuals}
\end{equation}
Let $\overline U$ consist of independent actual rows spanning the row
space of $U_0$, and let $\overline V$ consist of independent actual
columns spanning the column space of $V_0$. If these ranks are $s,t$, put
\begin{equation}
 K=\begin{pmatrix}0_{s\times t}&\overline U\\
                   \overline V&E\end{pmatrix}.
 \label{eq:core}
\end{equation}
Then
\begin{equation}
 \boxed{\rk(A+UV)=\rk A+\rk K-r.}
 \label{eq:rankformula}
\end{equation}
The core has $(s+r)\times(t+r)$ entries, with $s,t\le r$. The matrix $E$
may be singular or zero. No inverse of $E$ is used.
\end{theorem}
\begin{proof}
Consider the bordered matrix
\[
 B=\begin{pmatrix}A&U\\V&I_r\end{pmatrix}.
\]
Subtracting $U$ times the bottom block row from the top block row shows
that $\rk B=r+\rk(A+UV)$; subtraction equals addition in $\F$.

The reflexive identities decompose the domain of $A$ as
$\im h\oplus\ker A$ and its codomain as $\im A\oplus\ker h$. On the
first summands $A$ and $h$ are inverse isomorphisms. Choose compatible
bases so that both are the appropriate rectangular versions of
$\operatorname{diag}(I_q,0)$. In these bases write
$U=(U_1,U_2)^T$ and $V=(V_1,V_2)$. Eliminating the $I_q$ block from $B$
leaves
\[
 \begin{pmatrix}0&U_2\\V_2&I_r+V_1U_1\end{pmatrix},
\]
and therefore adds exactly $q$ to the rank. In the same bases $U_0$ has
rows $(0,U_2)^T$, $V_0$ has columns $(0,V_2)$, and $VhU=V_1U_1$.

It remains to compress the two zero-corner directions. An invertible row
operation on the top row block puts any chosen row basis of $U_0$ first
and makes all its other rows zero. Because the upper-left block is zero,
this does not change $E$ or $V_0$. The analogous column operation on the
left column block keeps a column basis of $V_0$ and zeroes the others,
without changing $E$ or the retained rows of $U_0$. Deleting the resulting
zero rows and columns leaves precisely \eqref{eq:core}, up to independent
basis changes in its top and left blocks. Thus $\rk B=q+\rk K$.
Equating the two expressions for $\rk B$ proves the result.
\end{proof}

\begin{example}[The smallest resonance]
For $A=U=V=(1)$ one has $A+UV=0$ and $E=0$. An inverse-based update is
undefined. Theorem~\ref{thm:border} instead has $s=t=0$, $K=(0)$, and
$1+0-1=0$ as required. This is the essential singular case, not an
exception to be handled by returning an inconclusive result.
\end{example}

\begin{corollary}[Exact homology from the core]
\label{cor:homology}
For a valid presentation \eqref{eq:model}, take
$h=\bigoplus_\nu h_\nu\otimes I$, using a reflexive inverse of every
template. Then
\begin{equation}
 \boxed{\beta(D)=\beta_0+2r-2\rk K.}
 \label{eq:homologycore}
\end{equation}
\end{corollary}
\begin{proof}
Substitute \eqref{eq:rankformula} into \eqref{eq:beta} and use
$\beta_0=N-2q$.
\end{proof}

The formula is independent of which rank minor produces $h$, although the
core matrix itself need not be canonical. It is also valid with redundant
ports and with zero-dimensional blocks. These edge cases are included in
the tests, rather than being excluded by an informal genericity assumption.

\subsection{Homological degrees can be retained}
The implementation reports total rank. The same theorem computes a graded
profile by applying it separately to each rectangular differential
$D_i:C^i\to C^{i+1}$, with the ports of degree $i$. If $R_i=\rk D_i$ and
$N_i=\dim C^i$, then
\[
 \dim H^i(C,D)=N_i-R_{i-1}-R_i.
\]
The total port count is the sum of these degree-group counts, so summing
the corresponding cubic bounds remains polynomial. This is an extension
of the proof, not a claim that the shipped command-line interface already
outputs quantum grading, integral torsion, representatives, or relative
chain maps.

\section{Finite-state multiplicity registers}
\label{sec:register}

\begin{definition}[Layered weighted register]
A length-$m$ register with $Q$ outputs consists of positive integers
$b_0,\ldots,b_m$, an initial row $\alpha\in\F^{b_0}$, two matrices
$T_j^0,T_j^1\in\F^{b_{j-1}\times b_j}$ for every layer, and a terminal
matrix $O\in\F^{b_m\times Q}$. It evaluates to
\begin{equation}
 F(x_1,\ldots,x_m)=\alpha T_1^{x_1}\cdots T_m^{x_m}O.
 \label{eq:register}
\end{equation}
Its width is $B=\max_j b_j$. All layers are explicitly listed; $m$ is not
an exponentially large repetition count encoded in binary.
\end{definition}

For block $\nu$ the output count is $Q=2d_\nu r$. The first $d_\nu r$
bits encode the rows of $U$ at one word, and the next $d_\nu r$ bits
encode the columns of $V$ at that word. This simultaneous encoding permits
shared finite-state structure among different port coordinates. It is not
a requirement to build $Q$ separate automata.

A product of literals is represented by a diagonal register: one state
per product term, each transition retaining or deleting that state according
to the selected literal value. In particular, a point indicator needs one
term, not the exponentially long expansion that an unfortunate polynomial
basis might produce. General dense transition matrices are also allowed.

\begin{theorem}[Reachable rows with actual witnesses]
\label{thm:reach}
The span of the $2^m$ evaluation rows $F(x)$ has a basis consisting of at
most $B$ actual evaluations. Such a basis and its witnessing words can be
computed using $O(mB^3+B^2Q)$ binary field operations. The witness
bookkeeping requires only $O(mB)$ predecessor records.
\end{theorem}
\begin{proof}
At layer zero retain $\alpha$ if it is nonzero. Suppose rows
$w_1,\ldots,w_s$, each realized by a prefix, span all reachable rows at
layer $j-1$. Every reachable row at layer $j$ lies in the span of
\[
 w_iT_j^0,\quad w_iT_j^1\qquad(1\le i\le s).
\]
Conversely these candidates are themselves realized by prefixes. Choose
an independent subset of the actual candidates, not merely arbitrary
row-reduced substitutes. There are at most $b_j\le B$ selected rows.
Induction proves both the spanning property and the existence of actual
prefix witnesses. At the terminal layer multiply by $O$ and select an
independent subset again.

There are at most $2B$ candidates per layer; multiplication and elementary
elimination cost $O(B^3)$ field operations. Terminal multiplication and
selection cost $O(B^2Q)$. Store each new witness as a predecessor pointer
and one appended bit. At most $B$ records are retained per layer. Recovering
the at most $B$ final words then costs $O(mB)$ steps. Copying whole prefixes
at every layer would instead introduce an avoidable quadratic factor in $m$.
\end{proof}

This is a specialization of the established forward-space idea for weighted
automata \cite{kiefer}. Keeping actual witnesses makes the implementation
and its audit particularly transparent: each claimed basis row is checked
by replaying an ordinary binary word.

\begin{theorem}[Exact parity pairings]
\label{thm:pair}
Let $w(x)=\alpha T_1^{x_1}\cdots T_m^{x_m}$ be the terminal state row,
and define
\begin{equation}
 W_0=\alpha^T\alpha,\qquad
 W_j=(T_j^0)^TW_{j-1}T_j^0+(T_j^1)^TW_{j-1}T_j^1.
 \label{eq:gram}
\end{equation}
Then $W_m=\sum_x w(x)^Tw(x)$. For terminal coefficient columns $a,b$,
\begin{equation}
 \sum_x (w(x)a)(w(x)b)=a^TW_mb.
 \label{eq:pair}
\end{equation}
All sums here are in $\F$.
\end{theorem}
\begin{proof}
Partition the length-$j$ words by their last symbol and apply induction
to the outer product of their state rows. The scalar formula follows by
multiplying by $a^T$ and $b$. Dense evaluation of the recurrence costs
$O(mB^3)$ field operations.
\end{proof}

\begin{remark}[A Gram matrix is not a rank test in characteristic two]
\label{rem:gram}
For the nonzero constant function on $\{0,1\}^m$, $m\ge1$, the evaluation
matrix has rank one while its self-pairing is $2^m=0$ in $\F$. Thus a Gram
matrix can vanish on a nonzero row space. The implementation uses reachable
spaces for ranks and the recurrence \eqref{eq:gram} only for contractions.
It never substitutes one task for the other.
\end{remark}

\subsection{How the core is computed without expansion}
Apply Theorem~\ref{thm:reach} to the simultaneous output register. For each
retained evaluation and each template coordinate, compute the corresponding
row of $(I-A_\nu h_\nu)U$ and column of $V(I-h_\nu A_\nu)$. The union of
these finite lists spans all rows and columns needed for
$\overline U,\overline V$. Indeed, each coordinate extraction and template
matrix action is linear in the simultaneous evaluation row.

The entries of $VhU$ are sums, over templates and words, of products of a
$V$ coordinate and an $hU$ coordinate. Theorem~\ref{thm:pair} computes these
sums by small terminal coefficient matrices and $W_m$. The implementation
does \emph{not} build a full $Q\times Q$ output Gram matrix. It contracts
the $B\times B$ state matrix directly against $B\times r$ port
coefficient arrays, one template coordinate at a time.

\subsection{Elementary closure properties and their cost}
A fixed linear transformation of the output functions changes $O$ only;
it leaves every transition matrix and the width unchanged. A sum of two
register functions has a direct-sum representation of width at most
$B_1+B_2$. Their pointwise product has a tensor-product representation of
width at most $B_1B_2$. These identities follow by distributing matrix
products or by the elementary identity
$(aA\otimes bB)=(a\otimes b)(A\otimes B)$.

These are exact operations over $\F$, with no numerical tolerance or
singular-value truncation. They also expose a danger: repeated pointwise
products can multiply widths. Closure under an operation is not a
polynomial bound on the size after a long sequence of such operations.

\section{Regular-language multiplicities and exact cardinalities}
\label{sec:language}

Boolean multiplicity domains are useful but need not include all words.
For example, a fixed homological layer may correspond to words of a fixed
Hamming weight. We allow an explicitly supplied \emph{deterministic} layered
language automaton, with at most $L$ states per layer, a single initial
state, and designated accepting terminal states.

The number $a=|\Omega|$ is obtained by dynamic programming with integer
counts: initialize the initial-state count to one, and for each symbol
add each state's count to the count of its specified successor. There
are $O(mL)$ additions of integers with at most $m+1$ bits. Determinism
matters: counting accepting paths of an arbitrary nondeterministic
presentation would not in general count accepted words.

\begin{proposition}[Language restriction without coordinate enumeration]
\label{prop:language}
A width-$B$ output register can be masked by the indicator of a supplied
width-$L$ deterministic language using a width-at-most-$BL$ register.
The core in Theorem~\ref{thm:border}, its validity checks, and the port
ranks can then be computed over the accepted domain without enumerating
accepted or rejected words. Exact dimension bookkeeping uses $|\Omega|$,
not its parity.
\end{proposition}
\begin{proof}
Take the tensor product of the weighted register with the deterministic
state-transition representation of the language. At the terminal layer
set an output row to zero unless the language state accepts. The value
on a word is precisely $\mathbf 1_\Omega(x)F(x)$. Hence the evaluation
row spaces are those from accepted words, together with extra zero rows.
The parity pairings are also correct because
$\mathbf 1_\Omega(x)^2=\mathbf 1_\Omega(x)$ in $\F$.

To see the dimension statement transparently, extend the masked model to
all Boolean words. Its differential is a direct sum of the desired
accepted-domain complex and
\[
 (T,A_T)^{\oplus(2^m-|\Omega|)}
\]
on the rejected domain. No port enters or leaves the latter summand.
Consequently its dimension, differential rank, and homology dimension
are respectively
\[
 d(2^m-|\Omega|),\qquad q_T(2^m-|\Omega|),\qquad
 (d-2q_T)(2^m-|\Omega|).
\]
Subtract these \emph{known exact integers} from the full-word calculation.
The small core is unchanged, since all rejected-domain port coordinates
are zero. This is also the implementation used by the optional language
wrapper. The argument applies blockwise even when the accepted blocks
are coupled to each other.
\end{proof}

\begin{remark}[Do not saturate before removing the known complement]
The full-word padded model can have enormous homology entirely supported
on rejected words. A capped value cannot be used for the exact subtraction
in Proposition~\ref{prop:language}. The implementation first removes the
known uncoupled complement using arbitrary-precision integers, and only
then forms decision caps for the actual domain. A test with a zero base,
$2^{40}$ padded copies, and only one accepted word has actual homology
zero despite the enormous padded complementary homology.
\end{remark}

The fixed-weight language $|x|=k$ has a deterministic counter with
$k+2$ states: counts $0,\ldots,k$ and a dead overflow state. Its accepted
cardinality is $\binom mk$. Thus the representation supports binomial,
not just power-of-two, multiplicities at polynomial descriptor cost.
No unproved statement about Khovanov diagrams having such a representation
is implicit in this observation.

\section{Checking the full differential and minimizing ports}
\label{sec:validity}

\subsection{A symbolic square-zero check}
It is insufficient to check the template differentials and assume their
coupling remains a complex. The following test evaluates all of $D^2$,
not a sample of its entries.

\begin{theorem}[Finite-state validity certificate]
\label{thm:square}
Assume $A^2=0$, and compute $M=VU$. Then
\begin{equation}
 D^2=\bigl[\,AU+UM\ \ U\,\bigr]
       \begin{bmatrix}V\\VA\end{bmatrix}.
 \label{eq:square}
\end{equation}
The condition $D^2=0$ is equivalent to every pairing between a basis of
the row space of the first factor and a basis of the column space of the
second factor being zero. Both bases have size at most $2r$ and are
constructible from the register evaluations without expanding $C$.
Homogeneity of all ports is checkable by the same evaluations.
\end{theorem}
\begin{proof}
Expanding $(A+UV)^2$ and using $A^2=0$ gives
$AUV+UVA+U(VU)V$, which is \eqref{eq:square}. Each entry of a product
matrix is a pairing of a row of the first factor with a column of the
second. Bilinearity makes vanishing on spanning bases equivalent to
vanishing on every such pair.

The parity-pairing algorithm computes $M$ exactly. At each template and
word, $AU+UM$ and $VA$ are linear expressions in the simultaneous
register output after $M$ is known. The reachable-row theorem therefore
provides spanning lists for all required rows and columns. Ordinary
elimination reduces them to at most $2r$ each.

A forbidden degree coordinate is likewise a fixed linear functional of
the simultaneous output. It vanishes on every word exactly when it
vanishes on a spanning set of evaluations. Checking these coordinates
therefore certifies port homogeneity, including functions that vanish
only after cancellation in the register. With language restrictions,
use the masked evaluations and the actual accepted domain.
\end{proof}

This check establishes an algebraic chain complex. It does not establish
that the complex comes from a planar diagram, nor that its coefficients
encode valid cobordisms before closure. Such geometry remains the
responsibility of the producer.

\subsection{The minimum number of ports for a fixed base}
Redundant ports waste cubic core work and weaken rank-budget obstructions.
Their redundancy can be removed exactly, without modifying register widths.

\begin{theorem}[Exact port minimization]
\label{thm:minimize}
For a fixed base differential $A$, the smallest number of homogeneous
scalar ports in a factorization of $D-A$ equals
\begin{equation}
 r_{\min}=\rk(D-A).
 \label{eq:rmin}
\end{equation}
For a port-register presentation this number and a factorization attaining
it are constructible in polynomial time in the presentation. Every
register retains its length, transition matrices, and widths; only its
terminal output coefficients change. This does not optimize the base $A$.
\end{theorem}
\begin{proof}
First work within one port-degree group and write the update as $UV$ with
$k$ ports. Choose an actual row basis $P\in\F^{a\times k}$ of $U$ and
an actual column basis $Q\in\F^{k\times b}$ of $V$. These are computed
using reachable register evaluations. Choose a right inverse $P^+$ and
a left inverse $Q^+$ using rank minors. Define
\[
 L=UP^+,\qquad R=Q^+V.
\]
Because every row of $U$ is in the row space of $P$, and every column of
$V$ is in the column space of $Q$, one has $U=LP$ and $V=QR$.
Furthermore, the actual rows of $U$ selected for $P$ give an identity
submatrix in $L$; the actual columns selected for $Q$ give one in $R$.
Hence $L$ has full column rank and $R$ has full row rank.

It follows that
\[
 \rk(UV)=\rk\bigl(L(PQ)R\bigr)=\rk(PQ).
\]
Factor the small matrix $PQ=ST$ using exactly $r'=\rk(PQ)$ columns in
$S$. Then
\begin{equation}
 U'=UP^+S,\qquad V'=TQ^+V
 \label{eq:portmin}
\end{equation}
satisfy $U'V'=UV$ and use $r'$ ports. Every entry of $U'$ and $V'$ is a
fixed linear combination of original port-coordinate functions. Thus
\eqref{eq:portmin} changes only terminal outputs of their registers.

Repeat separately for each degree. Distinct groups map distinct source
degrees to distinct target degrees, so the ranks of these blocks add.
The result has $\rk(D-A)$ homogeneous ports. Conversely any sum of $s$
rank-one channels has rank at most $s$, proving minimality. All matrices
used to form the port transformations have dimensions bounded by the
original port count. No basis vector of the virtual space is enumerated.
\end{proof}

\begin{example}[Duplicate channels can disappear completely]
Two identical columns of $U$ paired with two identical rows of $V$
contribute twice the same rank-one matrix and cancel over $\F$.
Both factors may be nonzero, but the optimal port count is zero. The
prototype checks this on a length-40 register. Minimizing $\rk U$ and
$\rk V$ separately would not detect every such cancellation: the small
interaction $PQ$ is essential. Nor can one substitute $\rk(VU)$ for
$\rk(UV)$: with an even number $M$ of coordinates, $U=\one$ and
$V=\one^T$ give $VU=0$ but $UV=J_M$ has rank one.
\end{example}

The theorem is a useful preprocessing operation because it improves both
the core dimension and the sharp rank budget. It is not an algorithm for
finding a low-rank perturbation of an arbitrary differential: the quality
of the \emph{chosen base} is a distinct and harder issue.

\section{The main presentation-size complexity theorem}
\label{sec:complexity}

Let the number of template blocks be $b$. For block $\nu$, let $d_\nu$ be
the template dimension, $m_\nu$ the explicitly listed register length, and
$B_\nu$ the maximum \emph{effective} width after any language mask. The
unmasked register width times the deterministic language width bounds
$B_\nu$. Define
\begin{equation}
 \begin{split}
 \Psi_\nu={}&d_\nu^3+(m_\nu+1)B_\nu^3+d_\nu^2B_\nu r\\
             &+d_\nu B_\nu^2r+d_\nu B_\nu r^2.
 \end{split}
 \label{eq:Psi}
\end{equation}

\begin{theorem}[Polynomial-time homology of supplied presentations]
\label{thm:main}
From an explicit port-register presentation one can deterministically
check the template and port degree conditions, check $D^2=0$, and, when
valid, compute $\beta(D)$ exactly without enumerating any $\Omega_\nu$.
The binary field-operation cost is
\begin{equation}
 O\!\left(\sum_{\nu=1}^{b}\Psi_\nu+r^3\right).
 \label{eq:maincost}
\end{equation}
The algorithm uses polynomial space in the descriptor. Integer
cardinality, index, and dimension arithmetic has polynomial bit cost in
the explicit input length. For supplied width-$L_\nu$ language automata,
a conservative bound for the cardinality additions is
$O(\sum_\nu(m_\nu+1)^2L_\nu)$ bit operations.
The optional exact minimization of Theorem~\ref{thm:minimize} preserves
this polynomial bound. No factor $2^{m_\nu}$ occurs in the field-operation
bound.
\end{theorem}
\begin{proof}
For each template, check $A_\nu^2=0$ and compute a reflexive inverse in
$O(d_\nu^3)$ field operations. Run the reachable-space and parity-pairing
algorithms. Their per-block work is bounded by
$O((m_\nu+1)B_\nu^3+d_\nu B_\nu^2r)$, since the simultaneous output
count is $2d_\nu r$.

Multiplying template matrices into the port coefficient arrays costs
$O(d_\nu^2B_\nu r)$. The contractions producing $VU$ and $VhU$ cost
$O(d_\nu B_\nu^2r+d_\nu B_\nu r^2)$. Acting on the at most $B_\nu$
retained evaluations to form quotient and validity rows has the same
bounds. Their total number is $O(\sum_\nu d_\nu B_\nu)$; each row has
at most $2r$ bits, so eliminating those lists costs
$O(\sum_\nu d_\nu B_\nu r^2)$. Checking all pairings of the two small
validity bases and reducing the core costs $O(r^3)$.

Equations \eqref{eq:basecounts} and \eqref{eq:homologycore} finish the
calculation. Full-word dimensions require only binary shifts. A language
count uses $O(m_\nu L_\nu)$ additions with $O(m_\nu)$-bit operands,
which gives the stated conservative cardinality cost. Dimensions and
ranks have at most
\[
 O\!\left(\max_\nu m_\nu+\log\bigl(1+\sum_\nu d_\nu\bigr)\right)
\]
bits. All explicitly supplied degree labels and matrix indices are also
accounted for by the input bit length.

The stored transition matrices, template matrices, terminal coefficient
arrays, small cores, and reachable witnesses occupy polynomial space.
In particular, witness construction uses shared predecessor records,
not arrays indexed by all words. Port minimization adds row-space
extraction and matrices of size at most $r$, together with fixed linear
output transformations; its cost fits the same polynomial estimate.
\end{proof}

The estimate is conservative elementary linear algebra, not an assertion
that Python takes one machine instruction per field operation. Arbitrary-
precision bit rows often make the implementation substantially faster
than that elementary estimate. Conversely the measured timings cannot
be used to replace the proof with a fitted exponent.

\subsection{A concrete algorithm contract}
A producer may provide a compressed presentation, but the consumer does
not trust cached ranks or an asserted square-zero identity. It performs
these steps:
\begin{enumerate}[label=\arabic*.,leftmargin=*,itemsep=2pt]
\item Validate template shapes and degrees, build any deterministic language
mask, and compute the accepted multiplicities as integers.
\item Compute template inverses, finite-state evaluation bases, parity
pairings, and the complete small square-zero check.
\item Optionally minimize the homogeneous ports while preserving the base.
\item Construct $\overline U,\overline V,E,K$, compute $\rk K$, and evaluate
\eqref{eq:homologycore}; remove any known padded complementary summands
before applying a decision cap.
\item Return the algebraic result and its replay data. Apply a topological
verdict only in a caller with independent diagram-to-complex provenance.
\end{enumerate}
The shipped \code{verify} command recomputes this contract and compares
its canonical output with the recorded certificate. It shares the solver
implementation and is therefore a \emph{replay checker}, not a separately
implemented proof checker. The independent small-matrix and crossing-cube
audits provide different checks, but are finite tests rather than formal
verification.

\subsection{What yields a quasi-polynomial recognition bound}
\begin{theorem}[Constructive, conditional recognition theorem]
\label{thm:conditional}
Let $n\ge2$ be the encoded input length of a validated knot diagram.
Suppose a proved, constructive procedure produces and certifies a
presentation whose homology is that of its closed unreduced Khovanov
complex over $\F$. Suppose its construction and provenance-verification
cost is $n^{O(1)}2^{O(\sigma(n))}$, the number of blocks and all register
lengths are polynomial in $n$, and
\[
 d_\nu,\ B_\nu,\ r\ \le n^{O(1)}2^{O(\sigma(n))}
\]
with the integer label lengths polynomially bounded as well. Then the
combined recognition procedure has bit complexity
\begin{equation}
 n^{O(1)}2^{O(\sigma(n))}.
 \label{eq:conditional}
\end{equation}
In particular $\sigma(n)=O(\log^2 n)$ gives $n^{O(\log n)}$ time.
\end{theorem}
\begin{proof}
Every term of \eqref{eq:maincost}, and the additional integer bookkeeping,
has the stated bound after substitution. Add the assumed production and
verification cost. The rank-two criterion from Section~\ref{sec:model}
then decides the input knot.
\end{proof}

This theorem is useful as an implementation target because the consumer
and all of its costs are now explicit. It is not a substitute for the
producer hypothesis. Merely postulating a small presentation, finding one
by exponential search, constructing one after an exponential scan, or
assuming a short certificate without an efficient verifier does not meet
the theorem. The largest unresolved task is controlling the construction
and maintenance of this representation on ordinary diagrams.

\section{Sharp rank budgets and safe decision saturation}
\label{sec:budget}

\begin{theorem}[Sharp perturbation bound]
\label{thm:budget}
For valid square-zero matrices $A,D$ on the same finite graded space,
with $D-A=UV$ and $r$ ports,
\begin{equation}
 |\beta(D)-\beta(A)|\le2\rk(UV)\le2r.
 \label{eq:budget}
\end{equation}
Consequently $\beta_0>2r+2$ is a rank-two obstruction. For the rank-two
query it is sufficient to retain the base count as
\begin{equation}
 \min(\beta_0,2r+3)
 \label{eq:cap}
\end{equation}
along with the core data. After port minimization the optimal fixed-base
budget is obtained by replacing $r$ with $r_{\min}$.
\end{theorem}
\begin{proof}
The inequalities
$\rk D\le\rk A+\rk(UV)$ and
$\rk A\le\rk D+\rk(UV)$ give the rank bound. Equation~\eqref{eq:beta}
turns it into \eqref{eq:budget}. If $\beta_0>2r+2$, the final homology
dimension exceeds two. Otherwise the capped value \eqref{eq:cap} retains
$\beta_0$ exactly, so \eqref{eq:homologycore} is still an exact test for
equality to two.
\end{proof}

\begin{proposition}[Sharpness and the failure of a fixed cap of three]
\label{prop:capsharp}
The obstruction threshold $2r+2$ cannot be lowered using the rank budget
alone. The direct-sum cap $\min(3,\beta_0)$ is insufficient for arbitrary
coupled presentations, even with one port.
\end{proposition}
\begin{proof}
Take $A=0$ on a two-term space of total dimension $2r+2$, with $r+1$
generators in each degree. Let $D$ contain $r$ disjoint identity arrows.
Then $D^2=0$, $\rk D=r$, and $\beta(D)=2$. Thus a base dimension of
$2r+2$ can still lead to rank two. Removing the two isolated generators
shows the full downward change $2r$ is attained. Conversely, starting
with $r$ contractible pairs and cancelling all their arrows by an equal
perturbation attains the upward change $2r$.

For the fixed-cap failure, take $r=1$. A zero base of dimension four
with a rank-one two-term differential has final homology dimension two.
Adding two isolated generators gives base dimension six and final
homology dimension four. In both examples the nontrivial core is the
same $2\times2$ matrix
\[
 \begin{pmatrix}0&1\\1&1\end{pmatrix}
\]
of rank two, and both base dimensions would be capped to three. The
capped base and unchanged core would no longer distinguish the cases.
\end{proof}

This is an extension-specific warning, not a defect in the maintained
literal-direct-sum sharing method. Nonnegative multiplicity saturation
is sound when no differential connects the summands. Coupling introduces
new cancellations, so the available cancellation rank must enter the
threshold.

\begin{corollary}[Necessary near-contractibility of a useful unknot base]
\label{cor:necessary}
If a fixed-base presentation represents a knot of total rank two, then
\[
 r_{\min}\ge \frac{\beta_0-2}{2}.
\]
In particular, a huge zero-differential base cannot be repaired into an
unknot complex using only a small number of scalar ports.
\end{corollary}

This gives a practical design constraint: a promising producer must
identify most contractible structure \emph{before} encoding the remaining
couplings. Assigning every expanded generator to a zero template and
hoping that a few global channels will reconstruct the differential is
not a viable general strategy.

\subsection{Future updates need a future budget}
Suppose, on a fixed ambient space, a current valid differential may later
receive updates whose ranks sum to at most $R$, and the final differential
is again valid. The same proof gives a possible homology change of at most
$2R$. A rank-two obstruction must therefore reserve that whole remaining
budget. A bound using only channels already processed would be unsound.
This observation does not automatically apply to future crossing additions,
which usually change the ambient complex and require relative gluing data.

\section{An obstruction to automatic scalar-port closure}
\label{sec:amplification}

Finite-state compression solves evaluation and rank extraction for the
specified model. It cannot by itself make the model closed under every
operation that occurs in a topological scanner.

\begin{theorem}[Tensor amplification of the minimum port count]
\label{thm:tensor}
Fix a base $A$ and an update $\Delta=D-A$ of rank $s>0$. If both the base
and differential are tensored with an identity on a $q$-dimensional
multiplicity space, then the minimum scalar-port count for that fixed
new base is exactly
\begin{equation}
 \rk\bigl(\Delta\otimes I_q\bigr)=sq.
 \label{eq:amplification}
\end{equation}
No choice of scalar register-coordinate functions can represent this
same fixed-base update with fewer than $sq$ ports.
\end{theorem}
\begin{proof}
Choose bases reducing $\Delta$ to an identity block of size $s$ and a
zero block. Tensoring these invertible basis changes with $I_q$ reduces
$\Delta\otimes I_q$ to an identity block of size $sq$ and a zero block.
The lower bound on any port factorization follows from rank. A rank
factorization realizes equality, so the bound is exact.
\end{proof}

Even appending $t$ completely passive binary multiplicity factors can
therefore change a nonzero fixed-base port requirement from $s$ to
$s2^t$. This is a genuine algebraic obstruction, not a poor choice of
weighted-automaton basis. It rules out the argument that a compact
register description alone guarantees small port count under arbitrary
tensor extension.

There are possible responses, but none is free: absorb the replicated
update into a newly chosen template; allow structured matrix-valued ports;
use a recursively factored coupling description; or prove that the
particular geometric operation has smaller effective rank after additional
certified cancellations. Each option changes the construction problem and
needs its own size theorem. The existing structured-tensor and radical
methods may help with such absorption, but their compatibility with this
closed scalar interface is not asserted without proof.

\section{Explicit families and what they demonstrate}
\label{sec:examples}

\subsection{Connected support with one port and rank-two homology}
Let $M=2^m$, $m\ge2$, and consider the two-term complex
\begin{equation}
 \F^M\xrightarrow{\ I_M+\one(e_0+e_1+e_2)^T\ }\F^M.
 \label{eq:connected}
\end{equation}
Use the two-dimensional template with one identity arrow, so the base
is contractible. The column $U$ is the constant function and the row $V$
is the sum of three point indicators. A diagonal register of width four
encodes them for every $m$.

\begin{proposition}[An exponentially large connected example]
\label{prop:connected}
The differential in \eqref{eq:connected} has rank $M-1$, and its total
homology dimension is two. Its undirected differential support graph is
connected. It has a one-port presentation with a $1\times1$ zero core,
independent of $m$.
\end{proposition}
\begin{proof}
Put $v=e_0+e_1+e_2$. Then $v^T\one=3=1$ in $\F$.
If $(I+\one v^T)z=0$, then $z=\one(v^Tz)$, so the kernel is contained
in the span of $\one$. That vector is in the kernel, giving nullity one
and rank $M-1$. The two-term cokernel has the same dimension.

For connectedness, the source vertices $0,1,2$ are each adjacent to all
targets except the corresponding target of the same index. The subgraph
on these six vertices is a six-cycle. Every target with index at least
three is adjacent to all three of these source vertices. Every remaining
source has its identity edge to its own target. Thus all vertices lie
in one component.

The template inverse gives $U_0=V_0=0$ and
$E=1+v^T\one=0$. The core consists of this scalar, so
\eqref{eq:homologycore} gives $0+2-0=2$.
\end{proof}

This is a strict representation distinction from sharing literal
connected components \emph{of the supplied matrix}: the input has no
nontrivial differential-component decomposition to share. It is not a
lower bound against every algorithm that may first change basis or
perform cancellation. In particular, explicit scalar elimination will
reduce this example once it has allocated or otherwise accessed the
relevant matrix entries. The advantage demonstrated here is that the
register algorithm need not allocate its virtual generators at all.

For $m=40$, the represented chain space has
$2^{41}=2,199,023,255,552$ generators. The program also checks $m=200$ and
$m=1000$ using the same one-port construction. These are \emph{abstract}
chain complexes. No knot diagrams with these prescribed presentations
were constructed in this work.

\subsection{Dense parity-sensitive coupling}
Replace $v$ by $\one$ and restrict the multiplicity domain to words of
Hamming weight $k$. Write $a=\binom mk$. The resulting nonzero block is
\begin{equation}
 I_a+J_a,\qquad J_a=\one\one^T.
 \label{eq:binomial}
\end{equation}
Only one rank-one channel is needed. The unmasked output register has
width one; the deterministic counter mask has width $k+2$.

\begin{proposition}[Binomial resonance]
\label{prop:binomial}
The two-term complex with differential \eqref{eq:binomial} has total
homology dimension
\begin{equation}
 \beta=\begin{cases}2,&\binom mk\text{ is odd},\\0,&\binom mk\text{ is even}.
 \end{cases}
 \label{eq:binomialbeta}
\end{equation}
The odd case occurs exactly when every nonzero binary digit of $k$ is
also a nonzero binary digit of $m$.
\end{proposition}
\begin{proof}
Here $E=1+a$ in $\F$. If $a$ is even, then $J_a^2=aJ_a=0$, so
$(I_a+J_a)^2=I_a$ and the differential is an isomorphism. If $a$ is odd,
the same kernel argument as above gives nullity one. The case $a=0$
has zero-dimensional chain groups and belongs to the even case.

For the parity criterion, write $m=\sum_jm_j2^j$. In $\F[t]$,
\[
 (1+t)^m=\prod_{j:m_j=1}(1+t^{2^j}).
\]
Each exponent in the product is the sum of a unique subset of the
indicated powers of two. The coefficient of $t^k$ is therefore one
exactly under the stated binary condition. This is the familiar binary
Lucas criterion; it is used here as an independent closed-form check,
not claimed as a new combinatorial theorem.
\end{proof}

The packaged language examples have the following exact values:
\begin{center}
\begin{tabular}{rrrr}
\toprule
$(m,k)$ & Multiplicity $a$ & Total dimension $2a$ & $\beta$\\
\midrule
$(63,31)$ & 916312070471295267 & 1832624140942590534 & 2\\
$(64,32)$ & 1832624140942590534 & 3665248281885181068 & 0\\
\bottomrule
\end{tabular}
\end{center}
The first case has singular interaction and the second has invertible
interaction. Both are processed by the same formula. Integer cardinality
is essential for the virtual dimension, while parity is the correct
arithmetic for $VhU$; conflating these two roles would destroy the result.

\subsection{Large surviving homology is also represented exactly}
Take a zero base and the rank-one two-term differential $\one\one^T$ on
$M$ generators in each degree. Its homology dimension is $2M-2$ for
$M\ge1$. The register algorithm outputs this integer without a basis of
that size. This example tests the opposite regime from the contractible
base: a small number of ports cannot remove the large base homology,
exactly as Theorem~\ref{thm:budget} predicts.

\section{Integration with the maintained recognizer}
\label{sec:integration}

\subsection{A deliberately narrow adapter}
The module \code{portkh.explicit} contains an explicit audit bridge. It
accepts a scalar, graded, square-zero matrix; chooses vertex-disjoint
arrows as a square-zero base; factors the residual separately in every
degree; and constructs a zero-length register. Equality of the resulting
matrix with the input is checkable by expansion on small cases.
This is a correctness bridge, not an asymptotic compression algorithm:
the entire explicit matrix is already present and the template can be
as large as that matrix.

The companion \code{from_closed_fastscan} adapter targets the inspected
interned-list API. It refuses a nonempty frontier, non-scalar coefficients,
inconsistent live counts, mixed closed matching identifiers, and edges
targeting deleted objects. It was exercised with API-shaped fixtures and
independent small cubes. It was \emph{not} exercised against a complete
upstream checkout, the maintained 277-test suite, or the production PD
corpus. These distinctions are also recorded in the integration notes.

Running this adapter after full scalar elimination is generally redundant:
the surviving closed complex already has zero differential. Running it
before that elimination still cannot recover the cost of a prefix scan
that has already expanded exponentially. A genuine improvement needs a
producer that emits shared templates and coordinate registers earlier,
with the required relative information preserved until closure.

\subsection{Data that must not be discarded at an open frontier}
Before closure, differential coefficients are typed cobordisms rather than
ordinary scalars. Surviving matching multiplicities do not determine all
maps; the maintained radical discussion already gives examples of this
issue \cite{repo-radical}. The present use of a generalized inverse over
$\F$ does not bypass that category structure. In particular, it does not
license scalar mixing of inequivalent boundary objects, nor transfer of
a common-disk radical exponent to an uncertified arbitrary PD frontier.

A future relative adapter must state the object types of its ports,
the coefficient algebra of its maps, and the continuation information
carried by the compressed basis. A theorem about the rank of one closed
complex does not imply equivalence under every future gluing operation.

\subsection{An adoption gate based on structural cost, not optimism}
The natural first integration is an opt-in diagnostic, not replacement of
\code{scan_fast.py}. A producer should report template dimensions, raw
and minimized port counts, effective register widths, accepted-word
counts, estimated descriptor storage, and the cost of obtaining all of
those data. It should preserve the current fallback state until the new
operation has passed its checks.

For ordinary already-expanded matrices, the bridge should remain disabled
by default: the measured controls below show pronounced overhead. For a
native compressed producer, an admission test can compare a bound such
as \eqref{eq:maincost} with an explicit-allocation estimate. Such a test
is a resource policy, not a knot invariant. A declined conversion or
exhausted budget must fall back or report unknown, never be interpreted
as proof of nontriviality.

The supplied package does not modify the repository, enable a new default,
or assert production compatibility beyond the inspected interface. The
recommended integration directory is a new named research subdirectory,
not an assumed unused numeric report identifier.

\section{Validation and measured performance}
\label{sec:experiments}

\subsection{What was actually run}
The standard-library test suite passes 48 tests. Several tests include
seeded loops: 1,000 random rectangular rank updates, 300 reflexive-inverse
checks, 200 rank factorizations, 300 random finite-state reachability and
pairing comparisons, 250 coupled two-term examples, 160 three-term
validity trials, 100 exact port-minimization comparisons, and 70
language-restricted dense comparisons. Additional tests cover literal
truth tables, long shared witnesses, malformed input, degree failures,
empty dimensions, duplicate ports, unsafe saturation, large dimensions,
CLI replay, and certificate tampering. These are not all independent
mathematical tests of the same strength; the raw test log identifies
exactly what each method checks.

The separate validation script exhausts all binary $2\times2$ bases with
one or two ports. This gives 4,352 update cases, including 1,728 with a
singular middle matrix. Every result agrees with a distinct low-pivot
rank implementation. The unit suite additionally compares against a
dense list-of-bits elimination routine, rather than relying exclusively
on the production bit-row elimination.

An independent braid-closure cube constructor audits every signed
2-strand word of length at most five and every signed 3-strand word of
length at most four: 404 distinct presentations, of which 210 have one
component and 194 have more than one. Many presentations represent the
same knot or link. The largest cube dimension is 246. Every cube is
checked for $d^2=0$, converted through the explicit bridge, re-expanded
for equality, and compared in total homology. Each also has a mirror
homological-profile check. The mirror set overlaps the original set;
it is not another 404 distinct knot types.

Known small examples and tests of braid relations, cancellation, and
positive and negative stabilization supplement this enumeration. The
cube oracle is independent of the register and port construction, but
is itself a new small implementation rather than an external production
knot package. These finite checks support the code; they do not replace
the proofs or constitute formal verification.

The files \code{data/validation.json}, \code{data/unit-tests.txt},
\code{data/benchmarks.json}, and the example certificates preserve the
actual results. The scripts contain fixed random seeds. No full upstream
checkout validation or end-to-end production recognition benchmark was
run in this session.

\subsection{Paired timings and their interpretation}
All reported timings use seven rounds with randomized order among a
baseline, the port implementation, and a second identical baseline.
The A/A ratio measures local timing variability. Small cases use batches
of five calls. Imports, parsing, and disk I/O are excluded. Host and
Python metadata, every order, and every raw duration are stored in the
JSON, so the displayed ratios can be reconstructed. A paired ratio is the
median of within-round baseline/treatment ratios, not necessarily the
quotient of the two marginal median times.

For the connected abstract family, the baseline is intentionally favorable:
it constructs only the nonzero $M\times M$ block from its closed formula,
then performs bit-packed elimination. It does not allocate a
$2M\times2M$ total differential. The treatment analyzes an already
prepared compact register, including degree and square-zero validation.
Both therefore start from compact descriptions, but their construction
and checking work is not identical; the exact timed tasks are recorded.
The family also has a closed-form answer: this benchmark compares two
representation paths, not the best possible specialized algorithm for
this family.

\begin{table}[htbp]
\centering
\caption{Connected abstract examples: prepared-register analysis versus construction and rank of the explicit nonzero block. Seven randomized paired rounds. Times are milliseconds.}
\label{tab:abstract}
\begin{tabular}{rrrrrr}
\toprule
$m$ & Virtual dimension & Explicit & Register & Paired ratio & A/A \\
\midrule
4 & 32 & 0.0021 & 0.0832 & 0.025 & 0.913 \\
6 & 128 & 0.0071 & 0.1025 & 0.073 & 0.944 \\
8 & 512 & 0.0281 & 0.1143 & 0.245 & 0.989 \\
10 & 2,048 & 0.1325 & 0.1369 & 0.962 & 1.039 \\
12 & 8,192 & 0.5673 & 0.1637 & 3.491 & 1.001 \\
14 & 32,768 & 7.0811 & 0.2394 & 31.631 & 1.090 \\
\bottomrule
\end{tabular}
\end{table}
\begin{table}[htbp]
\centering
\caption{Compressed-only connected examples. No explicit competitor was run at these sizes; no speedup ratio is inferred.}
\label{tab:huge}
\begin{tabular}{rrrrrr}
\toprule
$m$ & Virtual dimension & Width & Ports & Time (ms) & Candidate rows \\
\midrule
40 & $2^{41}$ & 4 & 1 & 0.4141 & 310 \\
200 & $2^{201}$ & 4 & 1 & 2.0672 & 1590 \\
1000 & $2^{1001}$ & 4 & 1 & 11.7555 & 7990 \\
\bottomrule
\end{tabular}
\end{table}
\begin{table}[htbp]
\centering
\caption{Negative controls on prepared small knot cubes. The treatment includes the explicit-to-port producer. Times are milliseconds; a ratio below one is a slowdown.}
\label{tab:cubes}
\begin{tabular}{lrrrrr}
\toprule
Presentation & Dimension & Ports & Direct & Bridge & Paired ratio \\
\midrule
Stabilized unknot & 18 & 8 & 0.0032 & 0.2852 & 0.0113 \\
Trefoil & 30 & 15 & 0.0064 & 0.6077 & 0.0108 \\
Figure-eight & 66 & 36 & 0.0220 & 2.0978 & 0.0105 \\
Relator unknot & 90 & 49 & 0.0351 & 3.5818 & 0.0098 \\
\bottomrule
\end{tabular}
\end{table}
\newcommand{\SlowFour}{39.5}
\newcommand{\GainTen}{0.962}
\newcommand{\GainTwelve}{3.49}
\newcommand{\GainFourteen}{31.63}


The crossover is visible rather than hidden. At $m=4$ the register method
is approximately $\SlowFour$ times slower. At $m=10$ the median paired
ratio is $\GainTen$, close to the observed crossover. At $m=12$ and
$m=14$ the paired ratios are $\GainTwelve$ and $\GainFourteen$. These
are abstract-kernel observations on this host, not a fitted asymptotic
law and not knot-recognition speedups. The largest A/A variation in
these tiny measurements also argues against excessive precision in
performance claims.

The compressed-only cases demonstrate that the implementation follows its
representation-size theorem: length 40, 200, and 1000 examples return the
proved rank-two homology while the virtual dimensions are $2^{41}$,
$2^{201}$, and $2^{1001}$. No explicit run was attempted at those sizes,
and no hypothetical elapsed time or speedup is attributed to one.

The four small knot-cube controls make the limitation concrete. Their
prepared-cube direct baseline is much faster than the combination of the
explicit-to-port producer and the symbolic solver. Ratios around
$0.010$ correspond to roughly two orders of magnitude of overhead.
This is not a reason to discard the representation theorem, but it is a
reason not to install the audit bridge as a default optimization.

\section{Further research questions and prioritized tasks}
\label{sec:research}

\subsection{Native producers before expansion}
The first empirical task is to instrument actual scan stages and record
candidate template decompositions before object allocation. Can a stage
emit registers directly, rather than discover repeated rows after
expansion? A useful experiment must include production cost, certificate
cost, failed attempts, and peak storage. A small port count measured only
after exponential preprocessing is not progress toward the target bound.

\subsection{Choosing the base, not just minimizing its ports}
Theorem~\ref{thm:minimize} solves the optimal port count for a fixed $A$.
The next mathematical problem is to choose a compact template base that
simultaneously has small homology, small update rank, and a short
coordinate description. Corollary~\ref{cor:necessary} supplies a necessary
condition for unknot presentations. Can suitable bases be obtained from
certified local cancellation families rather than from general matrix
optimization?

\subsection{Relative, typed ports at tangle frontiers}
Can ports be defined in the relevant additive cobordism category so that
continuation preserves a compressed representation? The required theorem
must identify which basis changes respect object types, which contractions
remain valid after gluing, and how coefficient size enters the cost.
Merely applying the scalar formula to a residue complex loses information.

\subsection{Absorbing tensor amplification}
Theorem~\ref{thm:tensor} proves that naive scalar-port propagation fails
under passive multiplicity amplification. Investigate structured
matrix-valued channels or recursive factorizations that can absorb
$\Delta\otimes I$ without $q$ separate scalar ports. The central question
is whether their singular interactions still admit exact small-core
elimination with a proved size bound.

\subsection{Graded domains and simultaneous coordinate minimization}
The deterministic-language extension represents exact Hamming-weight
slices, but a practical constructor may create more complicated degree
conditions. Determine which conditions admit small deterministic
presentations and whether output registers can be minimized jointly with
those conditions. Avoid implicitly replacing a small nondeterministic
object by an exponentially large deterministic one.

\subsection{A diagram-level constructive parameter theorem}
Find a geometric or combinatorial parameter whose bounded value guarantees
an efficiently obtainable template-and-port presentation, with explicit
bounds for $d_\nu,B_\nu,r$, and the number of stages. Theorem~\ref{thm:conditional}
then turns that theorem into a runtime result. The parameter must be
computable, certified, or supplied with the input; existentially optimal
scan orders alone do not satisfy the contract.

\subsection{Adversarial lower bounds for register width}
Port rank and automaton width are independent sources of cost. Construct
families whose updates have small rank but whose coordinate functions
require large widths for natural variable orders. Conversely, identify
geometric transformations that change the variable order without
expanding the complex. Such lower bounds would clarify which representation
failures are intrinsic and which are artifacts of a poor layout.

\subsection{Composing with the existing restricted algorithms}
The twist, graded-occupancy, extremal-window, and radical-transfer methods
can produce manageable local complexes in regimes where generic scanning
is expensive. Study whether their interfaces give structured low-rank
couplings between blocks, rather than multiplying object counts. Every
composition theorem must count producer cost and preserve the exact
hypotheses of each local result, especially common-disk assumptions and
window guard degrees.

\subsection{Independent verification and formalization}
The replay checker currently shares its arithmetic implementation with
the solver. A smaller independent checker could verify rank factorizations,
reachability closure, witness evaluations, small pairings, and the rank
minor identities. Formalization should begin with Theorem~\ref{thm:border}
and the rank-minimization theorem, then establish the automaton and
language semantics. Topological provenance is a separate formalization
layer, not an automatic consequence of binary matrix verification.

\subsection{Production acceptance criteria}
Before a default change, run the full maintained suite, production
PD examples, several scan orders, and realistic hard-knot families.
Measure failed conversion overhead and capped fallback behavior. Only
then investigate a Rust port or more elaborate dense kernels. The present
positive abstract cases and negative knot controls define an honest
starting point for that work rather than a substitute for it.

\section{Conclusion}
The main completed advance is an exact polynomial-in-description algorithm
for homology of finite-state, port-coupled repeated templates. It remains
valid at singular interactions, checks the whole differential, supports
regular-language multiplicities, and removes every redundant scalar port
for a fixed base without increasing register width. These are executable
and mathematically specified operations, not an appeal to an unspecified
compression oracle.

The limits are equally concrete. A useful unknot base must already be
nearly contractible relative to its port budget; tensoring can force
exponential port amplification; open-frontier continuation needs typed
relative data; and the available explicit bridge is slower on the measured
small knot cubes. The next breakthrough needed for this route is therefore
a constructive theorem about \emph{producing and maintaining} suitable
representations. The consumer-side algebra and its exact cost accounting
are provided here as a platform for proving such a theorem.

\appendix
\section{Independent small-cube construction}
\label{app:cube}
The audit cube represents a braid with $s$ strands and $n$ crossings by
$s(n+1)$ level vertices. Unaffected strands are joined vertically. At an
active crossing, a smoothing either joins the two strands vertically or
joins the two upper endpoints and the two lower endpoints. For positive
crossings the two choices are assigned cube coordinates zero and one in
that order; for negative crossings the order is reversed. Closure joins
the first and last level at each strand. A union--find computation produces
the smoothing circles.

Each circle is labeled by $1$ or $x$ in $\F[x]/(x^2)$. The cube differential
uses the standard Frobenius maps
\[
 m(1,1)=1,\quad m(1,x)=m(x,1)=x,\quad m(x,x)=0,
\]
\[
 \Delta(1)=1\otimes x+x\otimes1,\qquad \Delta(x)=x\otimes x.
\]
An edge of the cube changes the number of circles by one. Comparing the
circle partitions identifies whether it is a merge or a split and which
labels are unaffected. The homological degree is the number of one
resolutions minus the number of negative crossings. Sign assignments are
unnecessary over $\F$. The standard Frobenius identities make square faces
commute; their two paths cancel in characteristic two, giving $d^2=0$.
The script additionally computes and checks that square directly.

The oracle has strict crossing and dimension caps. It is deliberately
exponential and independent of the register representation. Its purpose is
to validate small closed complexes and the explicit bridge, not to provide
an alternative scalable recognizer. The supplied figure-eight example
records the exact braid word, coefficient field, grading convention,
constructor, and resulting total rank ten.

\section{File formats, replay, and reproducibility}
\label{app:files}
The matrix convention is fixed throughout: an integer row has its column
zero coefficient in the least significant bit. The JSON format
\code{portkh-1} stores template differential rows, template degrees, port
degrees, and the complete register matrices. At one template word, output
bits are grouped by template coordinate: first $d r$ bits for the rows of
$U$, followed by $d r$ bits for the columns of $V$. Thus a $V$ column is
stored as an $r$-bit word, avoiding ambiguity about matrix orientation.

A language descriptor supplies its layer widths, deterministic successor
arrays for both symbols, initial state, and accepting states. Language
files are arrays aligned with the template blocks; a null entry means
all words. The optional language wrapper masks ports and removes the
known uncoupled complement exactly, as proved in
Proposition~\ref{prop:language}.

From the package root, no dependency installation is required:
\begin{verbatim}
export PYTHONPATH="$PWD/src"
python -m unittest discover -s tests -v
python scripts/validate.py
python scripts/benchmark.py
python scripts/make_tables.py
sh build.sh
\end{verbatim}
The basic large certificate can be replayed by
\begin{verbatim}
python -m portkh verify \
  examples/connected_singular_40.json \
  examples/connected_singular_40.certificate.json
\end{verbatim}
For the language-restricted singular example:
\begin{verbatim}
python -m portkh verify \
  examples/binomial_63_31.json \
  examples/binomial_63_31.certificate.json \
  --languages examples/binomial_63_31.languages.json
\end{verbatim}
The \code{minimize} command writes a new basic presentation with the
minimum homogeneous port count for its fixed base. It does not accept
language files directly; a program using language restrictions should
apply the proven mask first and keep the accepted-domain bookkeeping.
No numerical randomization occurs in the solver. Randomness is used only
for reproducible tests and randomized benchmark order.

The archive includes the full source, example inputs and certificates,
raw validation and timing records, integration notes, an MIT-0 license,
a source-provenance manifest, and file checksums. The PDF and its TeX are
both included. Compiling the article uses the adjacent generated
\code{measurements.tex}. Re-running benchmarks changes timing-dependent
files; the reproduction script regenerates tables and checksums afterward.
The measured numbers in the text describe the packaged run, not a promise
that another host will reproduce identical timings.

\begin{thebibliography}{99}
\small\setlength{\itemsep}{2pt}\setlength{\parskip}{0pt}
\bibitem{barnatan}
D.~Bar--Natan, \emph{Fast Khovanov homology computations},
Journal of Knot Theory and Its Ramifications 16 (2007), 243--255.
\href{https://arxiv.org/abs/math/0606318}{arXiv:math/0606318}.

\bibitem{kelomaki}
T.~Kelom\"aki and D.~Sch\"utz, \emph{On computational complexity of Khovanov
homology}, 2026 preprint, version 1.
\href{https://arxiv.org/abs/2601.02119}{arXiv:2601.02119}.

\bibitem{kronheimer}
P.~B.~Kronheimer and T.~S.~Mrowka, \emph{Khovanov homology is an
unknot-detector}, Publications Math\'ematiques de l'IH\'ES 113 (2011).
\href{https://arxiv.org/abs/1005.4346}{arXiv:1005.4346}.

\bibitem{shumakovitch}
A.~Shumakovitch, \emph{Torsion of the Khovanov homology},
Fundamenta Mathematicae 225 (2014), 343--364.
\href{https://arxiv.org/abs/math/0405474}{arXiv:math/0405474v2}.
In particular, Lemma 3.1.C and Corollary 3.2.C.

\bibitem{kiefer}
S.~Kiefer, \emph{Notes on Equivalence and Minimization of Weighted Automata},
2020. \href{https://arxiv.org/abs/2009.01217}{arXiv:2009.01217}.

\bibitem{oseledets}
I.~V.~Oseledets, \emph{Tensor-Train Decomposition},
SIAM Journal on Scientific Computing 33 (2011), 2295--2317.
\href{https://doi.org/10.1137/090752286}{doi:10.1137/090752286}.

\bibitem{lackenby}
M.~Lackenby, \emph{Incompressible surfaces, hierarchies and unknot recognition},
25 July 2026, version 1, especially Section 9.
\href{https://arxiv.org/abs/2607.23350}{arXiv:2607.23350}.

\bibitem{repo-readme}
V.~Reshetnikov, ProveIt repository, UnknotRecognition README.
Inspected 8 October 2026; blob prefix \code{7d80dc0740b2}.
Full identifiers and inspection scope are in \code{provenance.json}.

\bibitem{repo-scan}
ProveIt, \path{Topology/UnknotRecognition/fast/fastunknot/scan_fast.py}.
Inspected blob prefix \code{6c731d0f4d2c}; verified at commit
\code{81a387aada56}.

\bibitem{repo-structural}
ProveIt, \path{Topology/UnknotRecognition/synthesis/structural.tex}.
Inspected blob prefix \code{26b056b9a6f8}.

\bibitem{repo-radical}
ProveIt, \path{Topology/UnknotRecognition/synthesis/radical.tex}.
Inspected blob prefix \code{a163430fd6cb}.
\end{thebibliography}
\end{document}
